% =====================================================================
\chapter{Levels as Units of Verification-Obligation Decomposition}\label{ch:levels}
% =====================================================================

The first seven chapters use levels as a tool for comprehension
(Chapter \ref{ch:comprehension}), for design (Chapter \ref{ch:design}) and for assurance
(Chapter \ref{ch:intent}). This chapter shows that the three are three uses of one object,
and argues for the direction among them most worth pursuing: promoting levels from an aid to
comprehension to the \emph{unit into which verification obligations decompose}.

\section{The Cognitive Scale}\label{sec:levels-def}

\begin{definition}[Cognitive scale]\label{def:cognitive-scale}
Given a study object $S$ and a task $T$, the \term{cognitive scale} of $S$,
$\Lk{0}, \Lk{1}, \Lk{2}, \dots$, is divided by the type of question a level answers rather
than by the physical location of the code:
\begin{description}
\item[$\Lk{0}$ Computational principles] What performs the computation? What constrains it?
      Memory, instruction execution, function calls, concurrency, synchronisation, virtual
      memory, I/O, caching, networking, persistence, types and the memory model.
\item[$\Lk{1}$ Infrastructure and runtime] Which shared capabilities does the system
      provide, and how is their correct cooperation guaranteed? Allocators, task schedulers,
      event loops, object systems, executors, storage-engine interfaces, resource managers,
      lifecycle mechanisms. This is the level at which the six functions of
      Chapter \ref{ch:infrastructure} reside.
\item[$\Lk{2}$ Components and subsystems] Which duties exist? Where are the boundaries? How
      do they cooperate? The transaction subsystem, the rendering engine, the
      entity--component system, the network protocol stack, the replication module, the query
      executor.
\item[$\Lk{3}$ Domain function and business semantics] What does the system promise
      externally? Which behaviour is correct? Query semantics, game rules, layout rules,
      transaction isolation semantics, permissions and business invariants.
\item[$\Lk{4}$ and above: scenario and system-level behaviour] How does the whole system
      behave in real scenarios? End-to-end workflows, failure recovery, system-level
      performance, deployment topology, business objectives.
\item[$\Lk{\infty}$ Purpose and responsibility] What should be wanted? Which trade-offs are
      acceptable? This level belongs to humans; the reason is given in \S\ref{sec:linf}.
\end{description}
\end{definition}

\begin{principle}[The cognitive scale is not a physical hierarchy]\label{pr:not-physical}
The levels of Definition \ref{def:cognitive-scale} denote \emph{cognitive scales}, not
necessarily the physical layers of the code. The $\Lk{3}$ domain function of a game engine
may call the $\Lk{1}$ runtime interface directly; the transaction semantics of a database may
equally impose constraints on the storage layer.
\end{principle}

\begin{principle}[Relativity and recursiveness of levels]\label{pr:relative}
The same component may occupy different levels under different study tasks. An asynchronous
runtime is infrastructure ($\Lk{1}$) to an ordinary application developer, and to its own
developers a complex system requiring deep modification (it has a complete
$\Lk{0} \dots \Lk{4}$ of its own). Levels are therefore defined relative to the study object
and the current task.
\end{principle}

Principle \ref{pr:relative} has an important corollary: \emph{every large system contains a
complete level ladder internally}, and this embedded ladder is a distinct thing from the
external technology-stack ladder. This explains directly the conclusion of
Chapter \ref{ch:infrastructure}: every large system grows a ``small operating system'' inside
its own domain, and that is precisely its embedded $\Lk{1}$.

\begin{remark}[Difference from the causal hierarchy]\label{rmk:pearl}
The levels of this paper resemble the causal hierarchy \cite{pearl2009} in form: an
intervention at a lower level does not automatically transmit to a higher level, just as a
modification of the implementation layer does not automatically become a modification of the
contract layer. The two differ, however. The causal hierarchy distinguishes \emph{types of
query} (observation, intervention, counterfactual); the levels of this paper distinguish
\emph{observable domain and guarantee obligation} (Definition \ref{def:level}). Conflating
them yields a concrete error: the belief that being able to ask a question at a higher level
amounts to being able to intervene at that level, when the latter requires an interface the
former does not provide.
\end{remark}

\section{The Level Contract}\label{sec:four-tuple}

\begin{definition}[The level-contract quadruple]\label{def:level-contract}
A level is recognised only if one can write down the quadruple
\begin{equation}\label{eq:level-contract}
\mathrm{LC}(\Lambda) \;=\;
\bigl(\text{participants},\; \text{guaranteed invariants},\; \text{contract form},\;
\text{verification method}\bigr),
\end{equation}
together with the fifth item this paper additionally requires, the \term{oblivious set}
$\Obl(\Lambda)$ (Definition \ref{def:oblivious}), that is, what the level is permitted never
to care about.
\end{definition}

\begin{principle}[Criterion of level reality]\label{pr:level-real}
A level is real if and only if the five items of Definition \ref{def:level-contract} can be
written down. A ``level'' whose fifth item (the oblivious set) cannot be written down is not
a level but a directory or a namespace.
\end{principle}

Applying Definition \ref{def:level-contract} to the instances of
Table \ref{tab:level-instances} shows the concrete form taken by one and the same set of
questions at different levels.

\begin{table}[htbp]
\centering
\footnotesize
\caption{Instantiation of the level-contract quadruple on a relational database
engine}\label{tab:level-instances}
\begin{tabular}{L{1.2cm}L{2.4cm}L{3.6cm}L{3.4cm}L{3.6cm}}
\toprule
\textbf{Level} & \textbf{Participants} & \textbf{Guaranteed invariants} &
  \textbf{Contract form} & \textbf{May ignore} \\
\midrule
$\Lk{0}$ & pages, cache lines, instructions &
  the locality principle; atomicity of a cache line &
  hardware timing and the memory model &
  upper-level semantics; the notion of a transaction does not exist \\
$\Lk{1}$ & buffer pages, log records, locks, memory contexts &
  a dirty page is not reused before it is flushed; the log reaches disk before the data &
  the \code{palloc} contract, the buffer pin/release protocol &
  SQL semantics; query-plan shape \\
$\Lk{2}$ & parser, planner, executor, access methods &
  the tuples consumed by the executor satisfy the visibility of the current snapshot &
  the executor interface (open/next/close) &
  the physical layout of storage; the log format \\
$\Lk{3}$ & transactions, statements, sessions &
  the visibility and atomicity declared by the isolation level &
  SQL semantics and the definition of isolation levels &
  buffer-management policy; index structures \\
$\Lk{4}$ & applications, operations, clients &
  an end-to-end workflow remains recoverable under failure &
  service-level objectives and recovery-point objectives &
  the shape of the internal execution plan \\
\bottomrule
\end{tabular}
\end{table}

The most informative column of Table \ref{tab:level-instances} is the last. It makes explicit
the other half of ``what this level is responsible for'', namely ``what this level is
permitted not to know''. That column is the one most often missing from the design documents
of large systems, and the most valuable.

\section{Decomposition of Verification Obligations}\label{sec:obligation-unit}

This is the central claim of the chapter.

\begin{principle}[The three identities of a level]\label{pr:three-identities}
One and the same level structure carries three functions at once:
\begin{enumerate}[label=(\roman*)]
\item \term{unit of comprehension}: it fixes the lower bound on ``how much information is
      needed to understand a system'', namely the composite of the level quotients
      (Proposition \ref{prop:oblivious});
\item \term{unit of design}: it fixes the structure of ``the read-set of a typical change''
      (Theorem \ref{th:centralise});
\item \term{unit of verification obligation}: it fixes \emph{which properties can be assured
      mechanically, which must be monitored at runtime, and which must be left to the human
      division of responsibility}.
\end{enumerate}
Item (iii) is the part that existing discussions touch on but do not develop.
\end{principle}

A level is naturally a unit of verification-obligation decomposition because it supplies two
things at once: an \emph{observable domain} ($\Obs_i$, what the level can see) and an
\emph{invariant set} ($\Inv_i$, what the level must guarantee). The form of a verification
obligation is precisely to decide a property over some observation domain, so every
verification obligation attaches naturally to some level.

\begin{definition}[Level attribution of a verification obligation]\label{def:obligation-assign}
Given a property $\varphi$ and a level $\Lambda$, $\varphi$ is \term{attributed to} $\Lambda$
if $\varphi$ can be expressed entirely within the observable domain $\Obs_\Lambda$ of
$\Lambda$. The level of attribution of $\varphi$ is the \emph{lowest} level that can express
it.
\end{definition}

Three classes of obligation follow:

\begin{enumerate}[label=(\alph*)]
\item \term{Mechanically decidable at the level}: $\varphi$ is expressible within the
      observable domain of $\Lambda$ and admits an exhaustible decision procedure. Assigned
      to mechanical assurance.
\item \term{Monitorable at the level but undecidable}: $\varphi$ is expressible within the
      observable domain but can be decided only at execution time, over traces already
      observed. Assigned to runtime assurance, with the three states distinguished by
      Principle \ref{pr:three-states}.
\item \term{Not expressible at the level}: $\varphi$ requires cross-level information or
      concerns purpose itself. Assigned to human attestation or consensus assurance.
\end{enumerate}

\begin{table}[htbp]
\centering
\footnotesize
\caption{Distribution of assurance types across levels (example)}\label{tab:assurance-matrix}
\begin{tabular}{L{4.2cm}L{3.2cm}L{3.2cm}L{3.6cm}}
\toprule
\textbf{Property} & \textbf{Level of attribution} & \textbf{Assurance mode} &
  \textbf{Evidence type} \\
\midrule
memory safety & $\Lk{0}$--$\Lk{1}$ & mechanical (type system / ownership / sandboxing) &
  compile-time checking, verifier \\
no out-of-bounds access and bounded loops & $\Lk{0}$--$\Lk{1}$ &
  mechanical (bytecode verifier) & static proof at admission \\
no unauthorised I/O & $\Lk{1}$--$\Lk{2}$ & mechanical (capability model / effect types) &
  compile-time $+$ runtime boundary \\
functional postconditions & $\Lk{2}$--$\Lk{3}$ &
  mechanical (Hoare / refinement types / contracts) & proof terms, SMT decision \\
resource bounds & $\Lk{1}$--$\Lk{2}$ & mechanical (resource analysis) $+$
  environmental assumptions & symbolic bounds, worst-case execution time analysis \\
real-time deadlines & $\Lk{2}$--$\Lk{4}$ & mechanical (within the hardware model) $+$
  runtime & timing analysis, monitoring \\
isolation-level semantics & $\Lk{3}$ & mechanical (model checking / contracts) &
  model-checking results \\
end-to-end availability & $\Lk{4}$ & monitoring $+$ consensus &
  operational metrics, service-level objectives \\
the specification is faithful to the intent & $\Lk{\infty}$ & consensus &
  adversarial review, attestation records \\
\bottomrule
\end{tabular}
\end{table}

Table \ref{tab:assurance-matrix} concretises the claim of this chapter. For a property, three
questions are asked in order: to which level does it belong; which assurance can that level
provide; and what are the assumptions of that assurance. These three questions replace
questions that admit no answer, such as ``do we need formal verification''.

\section{The Bearer of Responsibility}\label{sec:linf}

\begin{principle}[Responsibility cannot be transferred]\label{pr:responsibility}
Deciding what is wanted and bearing the consequences of that decision are two facets of one
act of agency. Hence $\Lk{\infty}$ does not belong to machines, not because machines are not
intelligent enough, but because responsibility must be borne by a subject capable of bearing
consequences.
\end{principle}

This principle must be made precise by the results of \S\ref{sec:consensus}. The residual
responsibility of humans is \emph{not individual but collective}: not ``one person reviews
the specification'' but ``a cognitive network reaches consensus on the specification through
adversarial debate''. That network may include human experts of different backgrounds, models
of different training provenance, and explicitly designed adversarial roles.

\begin{principle}[Diversity is irreplaceable]\label{pr:human-diversity}
What makes humans irreplaceable in the age of AI is not that individual human judgement is
better than that of an AI, but that humans are an irreducible \emph{source of diversity} in
the cognitive network. By Corollary \ref{cor:diversity}, the reliability of consensus is
governed by the independence of error modes, and the ultimate source of independence is
\emph{different histories of problem formation}. The value of humans therefore lies in the
diversity of their perspectives, not in the precision of their computation.
\end{principle}

\section{Level-Language Mismatch}\label{sec:antipattern}

\begin{principle}[Consistency of level language]\label{pr:level-language}
The gravest error in reading code is to understand the current level in the language of the
wrong level.
\end{principle}

\begin{example}[Concrete forms]\label{ex:wrong-language}
\begin{enumerate}[label=(\alph*)]
\item Understanding business validation logic by pointer analysis, treating ``is this
      transaction permitted'' as a question of memory aliasing.
\item Understanding buffer management in business semantics, looking for commercial rules in
      a page-replacement policy.
\item Understanding database isolation levels in the language of scheduler fairness: both are
      concurrency problems, but the correctness criterion of the former is a throughput and
      latency distribution, and that of the latter is visibility and atomicity.
\item Solving multi-tenant isolation with a type system, because isolation is a cross-cutting
      concern (Principle \ref{pr:crosscutting}) and a type system can express only part of it.
\end{enumerate}
\end{example}

The antipattern of Example \ref{ex:wrong-language}(d) has a more serious form: it is not
merely a wrong language but an attempt to press a cross-cutting concern into the level
structure. By Principle \ref{pr:crosscutting}, such an attempt is structurally doomed to fail,
and its failure is often misdiagnosed as ``layering that is not thorough enough''.

The empirical content of this chapter is the reproducibility of
Definition \ref{def:obligation-assign}: whether different analysts give the same level
attribution to the same property, and whether the assignment of
Table \ref{tab:assurance-matrix} reduces the omission of properties that should be assured
mechanically but are left to manual review. With a falling omission rate as the metric,
the claim is testable.
